\documentclass[../../main.tex]{subfiles}
\begin{document}

\begin{center}
  {\fontsize{48pt}{58pt}\selectfont \textbf{1999}\normalsize\textbf{年度}}
\end{center}

\begin{table}[htb]
  \centering
  \caption{出題テーマと難易度}
  \begin{tabular}{|c|p{6.2cm}|c|c|c|}
    \hline
    \textbf{問題番号} & \textbf{テーマ} & \textbf{思考量} & \textbf{計算量} & \textbf{総合} \\
    \hline
    第1問 & 三角関数の定義と加法定理 &  &  &  \\
    \hline
    第2問 & 複素数の漸化式と個数の極限 & 難 & 難 & 難 \\
    \hline
    第3問 & 四面体の辺を電流が通る確率 & 普 & 普 & 普 \\
    \hline
    第4問 & 直交する平面上の円板の半径 & 難 & 普 & 普 \\
    \hline
    第5問 & 二項係数の偶奇 & 難 & 難 & 難 \\
    \hline
    第6問 & 指数関数を含む定積分の下界 & 難 & 普 & 普 \\
    \hline
  \end{tabular}
\end{table}

\begin{table}[htb]
  \centering
  \caption{解答}
  \renewcommand{\arraystretch}{1.4}
  \begin{tabular}{|c|c|p{9cm}|}
    \hline
    \textbf{問題番号} & \textbf{小問} & \textbf{答え} \\
    \hline
    第1問 & (1) & 一般角の三角関数を単位円上の点の座標で定義する \\
    \cline{2-3}
    & (2) & 正弦・余弦の加法定理を証明する（証明問題） \\
    \hline
    第2問 & (1) & $\dfrac{3\cdot5^{n-1}}{4}<|z_n|<
    \dfrac{5^n}{4}$（証明問題） \\
    \cline{2-3}
    & (2) & $\displaystyle\lim_{r\to\infty}
    \dfrac{f(r)}{\log r}=\dfrac{1}{\log5}$ \\
    \hline
    第3問 & (1) & $p(1+2p-7p^3+7p^4-2p^5)$ \\
    \cline{2-3}
    & (2) & $p^2(1+2p-7p^3+7p^4-2p^5)^2$ \\
    \hline
    第4問 & -- & 半径の和の最大値は$\dfrac{9}{8}$ \\
    \hline
    第5問 & (1) & $m=2^k$なら$0<n<m$に対し${}_mC_n$は偶数
    （証明問題） \\
    \cline{2-3}
    & (2) & $m=2^k-1$（$k\in\mathbb{N}$） \\
    \hline
    第6問 & -- & $\displaystyle\int_0^\pi e^x\sin^2x\,dx>8$
    （証明問題） \\
    \hline
  \end{tabular}
\end{table}

\paragraph{試験形態}
6問．試験時間は要確認．

\end{document}
